\documentclass{article}
\usepackage[utf8]{inputenc}
\usepackage[T1]{fontenc}
\usepackage{geometry}
\usepackage{graphicx}
\usepackage{float}
\usepackage{hyperref}
\usepackage[english,indonesian]{babel}

\title{Optimasi Pemrosesan Template LaTeX Menggunakan Pendekatan Berbasis AI}
\author{Academic Support Agent}
\date{\today}

\begin{document}
\maketitle

\begin{abstract}

\end{abstract}

\chapter{PENDAHULUAN}


\chapter{TINJAUAN PUSTAKA}


% Original content (commented out):
% %%%%%%%%%%%%%%
% Modelo Projeto TCC I
% Autor: Leonardo Garcez Dalenogari Alba
% Ano: 2020
%%%%%%%%%%%%%%

\documentclass[
    12pt,
    oneside,
    openany,
    a4paper,
    english,
    brazil
]{article}

\usepackage{mathptmx}
\usepackage[T1]{fontenc}
\usepackage[utf8]{inputenc}
\usepackage[brazil]{babel}
\usepackage{indentfirst}
\usepackage{color}
\usepackage{graphicx,graphicx}
\usepackage{epsfig,subfig}
\usepackage{microtype}
\usepackage{imakeidx}
\usepackage{ragged2e}

\usepackage{lipsum}
\usepackage{amsmath,amssymb,mathrsfs}
\usepackage{setspace}
\usepackage{verbatim}
\usepackage{float}
\usepackage{tabularx}
\usepackage{multirow}

\usepackage{url}
\usepackage[table]{xcolor}

\usepackage{fancyhdr}
\pagestyle{fancyplain}
\fancyhf{}
\fancyhead[R]{\small\thepage}
\cfoot{}

\renewcommand{\headwidth}{16cm}
\renewcommand{\headrulewidth}{0pt} %zera linha de cabeçalho
\renewcommand{\footrulewidth}{0pt} %zera linha de rodapé

\usepackage{titlesec}
\usepackage{textcase,relsize}
%\def\myFBRSC#1{{\larger[0]{\MakeTextUppercase{#1}}}\scshape }
%\DeclareRobustCommand{\FirstBigRestSmallCaps}[1]{\myFBRSC #1}

\titleformat{\section}%
{\normalfont\scshape\Large\bfseries}{\thesection}{0.2em}{}

\titleformat{\subsection}%
{\normalfont\large\bfseries}{\thesubsection}{0.4em}{}

\titlespacing* % starred version: first paragraph is not indented
    {\section} % <command>
    {0cm} % <left>
    {0.6em} % <before-sep>
    {0.6em} % <after-sep>

\usepackage[abnt-emphasize=bf,alf]{abntex2cite}

\usepackage[left=3cm,top=3cm,right=2cm,bottom=2cm]{geometry}

\usepackage[nottoc,numbib]{tocbibind}

\usepackage{quoting}
\newenvironment{citacao}
    {
        \setlength\parindent{0cm}
        \fontsize{10pt}{12pt}\selectfont
        \begin{quoting}[rightmargin=0cm,leftmargin=4cm]
    }
    {
        \end{quoting}
    }

\usepackage{titletoc}
\addto\captionsbrazil{\renewcommand{\contentsname}{\hfill Sumário\hfill\vspace{1cm}}}

\titlecontents{section}
[0cm]             % left margin
{}             % above code
{\thecontentslabel\hspace{0.2em}}             % numbered format
{}             % unnumbered format
{ \titlerule*[1pc]{. . . . .}\contentspage \nobreak}             % filler-page-format, e.g dots

\titlecontents{subsection}
[0cm]
{}
{\thecontentslabel\hspace{0.2em}}
{}
{\titlerule*[1pc]{. . . . .}\contentspage \nobreak}

\include{configuracao/dados}

\begin{document}

\frenchspacing
\pagenumbering{arabic}
\setcounter{page}{1}

%%%%%%
%   Parte pré-textual
%%%%%%
\singlespacing

\include{pretextual/capa}
\include{pretextual/folha-rosto}

\thispagestyle{empty}
\tableofcontents
\clearpage

%%%%%%
%   Parte textual
%%%%%%

\renewcommand\thesection{\arabic{section}\hspace{0.2em}-{}}
\renewcommand\thesubsection{\arabic{section}.\arabic{subsection}}

\justify
\onehalfspacing
\setlength\parindent{1.5cm}
\setlength{\parskip}{0.6em}

\begingroup
\let\clearpage\relax
\section[Introdução]{Introdução}
\input{caps/introducao}
\endgroup

\begingroup
\let\clearpage\relax
\section[Justificativa]{Justificativa}
\input{caps/justificativa}
\endgroup

\begingroup
\let\clearpage\relax
\section[Problema a ser resolvido]{Problema a ser resolvido}
\input{caps/Problema}
\endgroup

\begingroup
\let\clearpage\relax
\section[Objetivos]{Objetivos}
\input{caps/Objetivos}
\endgroup

\begingroup
\let\clearpage\relax
\section[Fundamentação Teórica]{Fundamentação Teórica}
\input{caps/fundamentacaoTeoria}
\endgroup

\begingroup
\let\clearpage\relax
\section[Metodologia]{Metodologia}
\input{caps/metodologia}
\endgroup

\begingroup
\let\clearpage\relax
\section[Cronograma]{Cronograma}
\input{caps/cronograma}
\endgroup

%%%%%%
%   Parte pós-textual
%%%%%%

\singlespacing

\begin{flushleft}
    \begingroup
    \let\clearpage\relax
    \bibliography{bibliografia}
    \endgroup
\end{flushleft}

% Assinaturas
\vspace{3cm}
\begin{flushright}
    \rule{6cm}{1pt} \\
    \autor

    \rule{6cm}{1pt} \\
    \orientador
    
    %Comente as linhas abaixo se não tiver coorientador
    \rule{6cm}{1pt} \\
    \coorientador
\end{flushright}

\end{document}

\end{document}